
\documentclass{article}
\usepackage{colortbl}
\usepackage{makecell}
\usepackage{multirow}
\usepackage{supertabular}

\begin{document}

\newcounter{utterance}

\centering \large Interaction Transcript for game `eqbench', experiment `eq\_bench', episode 138 with reflexion{-}llama{-}3.1{-}8b{-}instruct.
\vspace{24pt}

{ \footnotesize  \setcounter{utterance}{1}
\setlength{\tabcolsep}{0pt}
\begin{supertabular}{c@{$\;$}|p{.15\linewidth}@{}p{.15\linewidth}p{.15\linewidth}p{.15\linewidth}p{.15\linewidth}p{.15\linewidth}}
    \# & \multicolumn{2}{c}{Player} && \multicolumn{2}{c}{Game Master} \\
    \hline

    \theutterance \stepcounter{utterance}  
    & & & \multicolumn{4}{p{0.6\linewidth}}{
        \cellcolor[rgb]{0.9,0.9,0.9}{
            \makecell[{{p{\linewidth}}}]{
                \texttt{\tiny{[P1$\langle$GM]}}
                \texttt{You have attempted this same task before and did not succeed. These are the plans you wrote after those attempts.} \\
\texttt{Trial \#1: New plan: I overestimated the likelihood of Henrietta feeling offended and indignant, and underestimated the likelihood of her feeling satisfied. Upon re{-}reading the dialogue, I realize that Henrietta's tone is polite and reserved, and she doesn't seem to be reacting strongly to Barty's suggestion. In fact, she seems to be politely dismissing his offer, which suggests that she's not feeling offended or indignant. On the other hand, she's not showing any enthusiasm or satisfaction with Barty's suggestion, which suggests that she's not feeling satisfied either. However, she does seem to be slightly annoyed or frustrated with Barty's persistence, which suggests that she might be feeling a bit anxious or irritated. Therefore, my revised scores are:} \\
\\ 
\texttt{Offended: 2} \\
\texttt{Satisfied: 0} \\
\texttt{Anxious: 6} \\
\texttt{Indignant: 3} \\
\\ 
\texttt{Trial \#2: New plan: I overestimated the likelihood of Henrietta feeling anxious and indignant, and underestimated the likelihood of her feeling satisfied. Upon re{-}reading the dialogue, I realize that Henrietta's tone is polite and reserved, and she doesn't seem to be reacting strongly to Barty's suggestion. In fact, she seems to be politely dismissing his offer, which suggests that she's not feeling offended or indignant. On the other hand, she's not showing any enthusiasm or satisfaction with Barty's suggestion, which suggests that she's not feeling satisfied either. However, she does seem to be slightly annoyed or frustrated with Barty's persistence, which suggests that she might be feeling a bit anxious or irritated. } \\
\\ 
\texttt{To improve my strategy, I will pay closer attention to Henrietta's tone and language, and try to identify the subtle cues that suggest her emotional state. I will also consider the context of the dialogue and the relationship between the characters, as these can influence the emotional dynamics of the conversation. } \\
\\ 
\texttt{Specifically, I will focus on the following:} \\
\\ 
\texttt{* Paying attention to Henrietta's language and tone, and trying to identify the subtle cues that suggest her emotional state.} \\
\texttt{* Considering the context of the dialogue and the relationship between the characters, and how these might influence the emotional dynamics of the conversation.} \\
\texttt{* Avoiding overestimating the likelihood of certain emotions, and instead trying to make more nuanced and context{-}specific predictions.} \\
\\ 
\texttt{By following this plan, I hope to improve my ability to predict the emotional responses of characters in dialogue, and to provide more accurate and insightful answers.} \\
\\ 
\texttt{Trial \#3: New plan: I overestimated the likelihood of Henrietta feeling anxious and indignant, and underestimated the likelihood of her feeling satisfied. Upon re{-}reading the dialogue, I realize that Henrietta's tone is polite and reserved, and she doesn't seem to be reacting strongly to Barty's suggestion. In fact, she seems to be politely dismissing his offer, which suggests that she's not feeling offended or indignant. On the other hand, she's not showing any enthusiasm or satisfaction with Barty's suggestion, which suggests that she's not feeling satisfied either. However, she does seem to be slightly annoyed or frustrated with Barty's persistence, which suggests that she might be feeling a bit anxious or irritated. } \\
\\ 
\texttt{To improve my strategy, I will pay closer attention to Henrietta's tone and language, and try to identify the subtle cues that suggest her emotional state. I will also consider the context of the dialogue and the relationship between the characters, as these can influence the emotional dynamics of the conversation. } \\
\\ 
\texttt{Specifically, I will focus on the following:} \\
\\ 
\texttt{* Paying attention to Henrietta's language and tone, and trying to identify the subtle cues that suggest her emotional state.} \\
\texttt{* Considering the context of the dialogue and the relationship between the characters, and how these might influence the emotional dynamics of the conversation.} \\
\texttt{* Avoiding overestimating the likelihood of certain emotions, and instead trying to make more nuanced and context{-}specific predictions.} \\
\\ 
\texttt{By following this plan, I hope to improve my ability to predict the emotional responses of characters in dialogue, and to provide more accurate and insightful answers.} \\
\\ 
\texttt{New plan: I will revise my scores to reflect a more nuanced understanding of Henrietta's emotional state. I will also pay closer attention to the context of the dialogue and the relationship between the characters, and try to identify the subtle cues that suggest Henrietta's emotional state.} \\
\\ 
\texttt{Offended: 1} \\
\texttt{Satisfied: 0} \\
\texttt{Anxious: 5} \\
\texttt{Indignant: 0} \\
\\ 
\\ 
\\ 
\texttt{Your task is to predict the likely emotional responses of a character in this dialogue:} \\
\\ 
\texttt{Barty: Ah, Mrs. Adams, how delightful to see you. I trust you are well?} \\
\texttt{Henrietta: Mr. Harrington, I'm as well as can be expected. What brings you here today?} \\
\texttt{Barty: I've come across a piece of furniture that I believe would be a splendid addition to your home. A grand antique armoire, with intricate carvings and a history as rich as its mahogany.} \\
\texttt{Henrietta: And why, pray tell, would I need an armoire, Mr. Harrington?} \\
\texttt{Barty: Well, Mrs. Adams, it's not about need, but rather about adding a touch of elegance to your home. It's a statement piece, a testament to your refined taste.} \\
\texttt{Henrietta: My home does not need any more statements, Mr. Harrington. It speaks volumes of my life as it is.} \\
\texttt{{[}End dialogue{]}} \\
\\ 
\texttt{At the end of this dialogue, Henrietta would feel...} \\
\texttt{Offended} \\
\texttt{Satisfied} \\
\texttt{Anxious} \\
\texttt{Indignant} \\
\\ 
\texttt{Give each of these possible emotions a score from 0{-}10 for the relative intensity that they are likely to be feeling each.} \\
\\ 
\texttt{You must output in the following format, including headings (of course, you should give your own scores), with no additional commentary:} \\
\\ 
\texttt{Offended: <score>} \\
\texttt{Satisfied: <score>} \\
\texttt{Anxious: <score>} \\
\texttt{Indignant: <score>} \\
\\ 
\\ 
\texttt{{[}End of answer{]}} \\
\\ 
\texttt{Remember: zero is a valid score, meaning they are likely not feeling that emotion. You must score at least one emotion > 0.} \\
\\ 
\texttt{Your answer:} \\
            }
        }
    }
     \\ \\

    \theutterance \stepcounter{utterance}  
    & \multicolumn{4}{p{0.6\linewidth}}{
        \cellcolor[rgb]{0.561,0.737,0.561}{
            \makecell[{{p{\linewidth}}}]{
                \texttt{\tiny{[P1$\rangle$GM]}}
                \texttt{Offended: 2} \\
\texttt{Satisfied: 0} \\
\texttt{Anxious: 4} \\
\texttt{Indignant: 0} \\
            }
        }
    }
    & & \\ \\

    \theutterance \stepcounter{utterance}  
    & & & \multicolumn{2}{p{0.3\linewidth}}{
        \cellcolor[rgb]{0.9,0.9,0.9}{
            \makecell[{{p{\linewidth}}}]{
                \texttt{\tiny{[GM$|$GM]}}
                \texttt{Offended: 2} \\
\texttt{Satisfied: 0} \\
\texttt{Anxious: 4} \\
\texttt{Indignant: 0} \\
            }
        }
    }
    & & \\ \\

    \theutterance \stepcounter{utterance}  
    & & & \multicolumn{2}{p{0.3\linewidth}}{
        \cellcolor[rgb]{0.9,0.9,0.9}{
            \makecell[{{p{\linewidth}}}]{
                \texttt{\tiny{[GM$|$GM]}}
                \texttt{Offended: 6} \\
\texttt{Satisfied: 2} \\
\texttt{Anxious: 0} \\
\texttt{Indignant: 8} \\
            }
        }
    }
    & & \\ \\

    \theutterance \stepcounter{utterance}  
    & & & \multicolumn{2}{p{0.3\linewidth}}{
        \cellcolor[rgb]{0.9,0.9,0.9}{
            \makecell[{{p{\linewidth}}}]{
                \texttt{\tiny{[GM$|$GM]}}
                \texttt{game\_result = LOSE} \\
            }
        }
    }
    & & \\ \\

    \theutterance \stepcounter{utterance}  
    & & \multicolumn{4}{p{0.6\linewidth}}{
        \cellcolor[rgb]{0.95,0.95,0.95}{
            \makecell[{{p{\linewidth}}}]{
                \texttt{\tiny{[P1$\leftrightarrow$P]}}
                \texttt{Post{-}episode Reflexion (not sent to GM):} \\
\texttt{New plan: I overestimated the likelihood of Henrietta feeling anxious and indignant, and underestimated the likelihood of her feeling satisfied. Upon re{-}reading the dialogue, I realize that Henrietta's tone is polite and reserved, and she doesn't seem to be reacting strongly to Barty's suggestion. In fact, she seems to be politely dismissing his offer, which suggests that she's not feeling offended or indignant. On the other hand, she's not showing any enthusiasm or satisfaction with Barty's suggestion, which suggests that she's not feeling satisfied either. However, she does seem to be slightly annoyed or frustrated with Barty's persistence, which suggests that she might be feeling a bit anxious or irritated.} \\
\\ 
\texttt{To improve my strategy, I will pay closer attention to Henrietta's tone and language, and try to identify the subtle cues that suggest her emotional state. I will also consider the context of the dialogue and the relationship between the characters, as these can influence the emotional dynamics of the conversation.} \\
\\ 
\texttt{Specifically, I will focus on the following:} \\
\\ 
\texttt{* Paying attention to Henrietta's language and tone, and trying to identify the subtle cues that suggest her emotional state.} \\
\texttt{* Considering the context of the dialogue and the relationship between the characters, and how these might influence the emotional dynamics of the conversation.} \\
\texttt{* Avoiding overestimating the likelihood of certain emotions, and instead trying to make more nuanced and context{-}specific predictions.} \\
\\ 
\texttt{By following this plan, I hope to improve my ability to predict the emotional responses of characters in dialogue, and to provide more accurate and insightful answers.} \\
\\ 
\texttt{Offended: 1} \\
\texttt{Satisfied: 0} \\
\texttt{Anxious: 5} \\
\texttt{Indignant: 0} \\
            }
        }
    }
    & \\ \\

\end{supertabular}
}

\end{document}
